\subsection{可逆的形式幂级数}

命题 5. --- 在单个未定元的形式幂级数环 A{[}{[}T{]}{]} 中，多项式 1 - T 是可逆的，并且我们有 \((1 - \mathrm{T})^{-1} = \sum_{n=0}^{\infty} \mathrm{T}^{n}\) 。

因为

\[
(1 - \mathrm{T}) \left(\sum_ {n = 0} ^ {\infty} \mathrm{T} ^ {n}\right) = \sum_ {n = 0} ^ {\infty} \mathrm{T} ^ {n} - \sum_ {n = 0} ^ {\infty} \mathrm{T} ^ {n + 1} = 1.
\]

命题 6. --- 设 \(u \in \mathbf{A}[[\mathbf{I}]]\) ；那么，\(u\) 在 \(\mathbf{A}[[\mathbf{T}]]\) 中可逆的必要且充分条件是它的常数项在 \(\mathbf{A}\) 中可逆。

假设存在 \(v \in A[[I]]\) 使得 \(uv = 1\) 。设 \(\alpha, \beta\) 分别为 \(u\) 与 \(v\) 的常数项，则 \(\alpha\beta = 1\) ，所以 \(\alpha\) 可逆。

反之，设 \(\alpha\) 为 \(u\) 的常数项，并假设它可逆。则存在形式幂级数 \(t \in \mathbf{A}[[\mathrm{I}]]\) 使得 \(u = \alpha(1 - t)\) 且 \(\omega(t) > 0\) 。现在，存在环同态 \(\varphi: \mathbf{A}[[\mathrm{T}]] \to \mathbf{A}[[\mathrm{I}]]\) 使得 \(\varphi(\mathrm{T}) = t\) ，而 \(1 - \mathrm{T}\) 在 \(\mathbf{A}[[\mathrm{T}]]\) 中可逆（命题 5）；因之 \(1 - t\) 在 \(\mathbf{A}[[\mathrm{I}]]\) 中可逆，从而 \(u\) 也可逆。

注记. --- 设 M 为所有常数项为 1 的形式幂级数构成的集合。由命题 6，M 在乘法下构成一个交换群；A{[}{[}I{]}{]} 的乘法群因而是 M 与 A 的乘法群的直积。我们将给 M 赋予由 A{[}{[}I{]}{]} 的拓扑诱导的拓扑。对每个 \(\beta \in \mathbf{N}^{(I)}\) ，我们已在 IV 第26页定义了理想 \(a_{\beta}\) （它是 A{[}{[}I{]}{]} 的理想）；则 \(1 + a_{\beta}\) 是 M 的一个子群，且族 \((1 + a_{\beta})\) 是 M 中 1 的一个基本邻域系。由于 M 中的乘法是连续的，我们看出 M 是一个拓扑群（《一般拓扑学》，III，第223页）；换言之，映射 \(f \mapsto f^{-1}\) 在 M 中是连续的。

设 K 为一个交换域，\(\mathfrak{D}\) 为有理分式域 \(\mathrm{K}((\mathbf{X}_i)_{i\in I})\) 中由使 \(\mathbf{K}^I\) 的元素 0 可代入的有理分式组成的子环。若 \(f\in \mathfrak{D}\) ，我们有 \(f = \frac{u}{v}\) ，其中 \(u\) 与 \(v\) 是多项式，且 \(v\) 的常数项 \(\neq 0\) ，因而 \(v\) 在 \(\mathbf{K}[[\mathrm{I}]]\) 中可逆。我们可以立即验证，元素 \(uv^{-1}\) （作为 \(\mathbf{K}[[\mathrm{I}]]\) 的元素）只依赖于 \(f\) ；我们称形式幂级数 \(uv^{-1}\) 为有理分式 \(\frac{u}{v}\) 在原点的展开式。映射 \(f\mapsto uv^{-1}\) 是从 \(\mathfrak{D}\) 到 \(\mathbf{K}[[\mathrm{I}]]\) 中的单同态；我们将经常把 \(\mathfrak{D}\) 与它在此映射下的像等同。

